\chapter*{Résumé et fiche du personnage}
\phantomsection
\addcontentsline{toc}{chapter}{Résumé et fiche du personnage}
\label{chap:resume-personnage}
\markboth{Résumé et fiche du personnage}{Résumé et fiche du personnage}

\section*{Fiche du personnage}
\phantomsection
\addcontentsline{toc}{section}{Fiche du personnage}
\label{sec:fiche-personnage}

\begin{description}
    \item[Nom :] Kael
    \item[Nom de naissance :] Khaer Ghal
    \item[Espèce :] Orc
    \item[Âge :] Entre 20 et 25 ans
    \item[Classe :] Rôdeur
    \item[Historique :] Voyageur
    \item[Alignement :] Chaotique Bon
    \item[Niveau :] 1
\end{description}

\illustration[0.95\textwidth]
  {fiche_personnage.png}
  {Fiche visuelle de Kael : apparence générale, équipement et détails caractéristiques.}
  {fiche-personnage-kael}

% À compléter après création de la fiche avec le MJ :
%
% Force :
% Dextérité :
% Constitution :
% Intelligence :
% Sagesse :
% Charisme :
%
% Classe d'armure :
% Points de vie :
% Bonus de maîtrise :
% Compétences :
% Langues :
% Équipement :
% Sorts :

\section*{Apparence}

Kael est un jeune Orc à la peau gris-bleu, aux yeux jaune doré et aux cheveux
sombres, courts et généralement en bataille. Son visage et sa carrure portent
les traits caractéristiques de son peuple, avec notamment deux défenses
inférieures visibles.

Son oreille gauche porte la marque de sa première épreuve : une créature lui en
a arraché une partie lorsqu'il était enfant, laissant le cartilage
irrégulièrement déchiré. Une longue canine prélevée sur cette même créature
traverse directement son oreille. La dent est portée brute, sans métal ni
ornement, seulement légèrement retaillée et entaillée afin de rester en place.

Autour du cou, Kael porte en permanence un pendentif transmis par sa mère lors
de sa fuite du village : une petite pierre sombre gravée d'un ancien symbole
orc. Il s'agit de l'un de ses biens les plus précieux.

\illustration[0.62\textwidth]
  {Kael_portrait.png}
  {Portrait de Kael.}
  {portrait-kael}

\section*{Personnalité}

Kael est profondément attaché à sa liberté. Il agit selon ce qu'il estime
juste et supporte mal qu'on décide à sa place de ce qu'il doit faire, de
l'endroit où il doit aller ou de la personne qu'il devrait devenir. Les règles,
les traditions et l'autorité n'ont à ses yeux de valeur que lorsqu'elles lui
paraissent avoir du sens.

Cela ne l'empêche pas de s'attacher profondément aux autres. Il tient à ceux
qu'il considère comme les siens et n'hésitera pas à les aider ou à les défendre
s'ils en ont besoin, mais cet attachement ne l'empêche pas de suivre ses propres
envies, de prendre ses propres décisions ou de partir seul lorsqu'il en ressent
le besoin. Il peut disparaître quelques heures pour explorer une forêt simplement
parce qu'il voulait savoir ce qui s'y trouvait, puis revenir auprès des siens
comme si cela allait de soi.

Sa morale est avant tout personnelle. Il supporte particulièrement mal
l'esclavage, la captivité, la violence gratuite et, plus généralement, tout ce
qui consiste à imposer par la force sa volonté à quelqu'un d'autre.

Il possède également un défaut particulièrement exploitable : son orgueil.
Mettre en doute ses capacités est souvent le meilleur moyen de le pousser à
essayer quelque chose qu'il aurait autrement jugé inutile ou dangereux. Un
simple « tu n'en es pas capable » ou « cap ou pas cap ? » peut parfois prendre
le dessus sur son bon sens.

Kael cherche avant tout à progresser. Il aime apprendre de nouvelles techniques,
observer ceux qui savent faire quelque chose qu'il ignore encore, essayer à son
tour et constater qu'il est devenu meilleur qu'auparavant. Il ne sait pas
exactement où cette recherche le conduira : il considère simplement qu'il doit
trouver et suivre sa propre voie.

\section*{Historique}

Khaer Ghal est né dans un village orc isolé dirigé par son père. Élevé selon
les traditions du clan, il apprend très jeune le combat, la chasse et la
survie. Lors de sa première épreuve, à huit ans, il tue seul une créature qui
lui arrache une partie de l'oreille gauche. L'une de ses canines devient son
premier trophée.

Quelques années plus tard, il se lie avec Kamatsu, un jeune humain retenu
prisonnier avec sa famille. Lorsque ceux-ci sont envoyés dans l'arène du
village, Khaer Ghal intervient pour les protéger. Après une première tentative
d'évasion qui coûte la vie aux parents de Kamatsu, sa propre mère aide
finalement les deux garçons à fuir.

Ils survivent seuls pendant une semaine avant de rencontrer une caravane qui,
après une période de méfiance, devient leur nouvelle famille. Khaer Ghal y
grandit avec Kamatsu et apprend au contact des nombreux voyageurs rencontrés
sur les routes : survie, chasse, combat, soins, observation, discrétion et,
plus tard, véritable maîtrise de l'arc.

C'est durant ces années que Kamatsu commence à l'appeler Kael, surnom que
Khaer Ghal finit par adopter lui-même. Leur amitié devient avec le temps un
véritable lien fraternel.

À l'âge adulte, leurs routes se séparent. Kamatsu poursuit sa propre vie tandis
que Kael finit par quitter à son tour la caravane. Il passe alors quelques
années à voyager seul, non pour atteindre une destination précise, mais pour
découvrir le monde, continuer à apprendre et trouver sa propre voie.

Ses voyages le conduisent finalement à Waterdeep où, sans beaucoup d'argent,
il tente de proposer ses services de rôdeur contre quelques pièces ou un
repas. Pris pour un mendiant par les gardes de la ville, il est arrêté et jeté
en prison.

C'est là que commence son aventure.

\illustration[0.48\textwidth]
  {kael_jeton.png}
  {Portrait de Kael utilisé comme jeton de personnage.}
  {jeton-kael}

\section*{Liens importants}

\begin{description}
    \item[Kamatsu :] Frère de cœur de Kael. Ils ont fui ensemble le village
    orc, grandi ensemble dans la caravane et se considèrent comme frères malgré
    l'absence de tout lien de sang.

    \item[Sa mère :] Elle a permis à Kael et Kamatsu de fuir le village et lui
    a transmis son pendentif avant leur départ. Kael ignore ce qu'elle est
    devenue.

    \item[La caravane :] Sa famille d'adoption et l'endroit où il a passé
    l'essentiel de son enfance et de son adolescence. Une grande partie de ce
    qu'il sait provient des voyageurs qu'il y a rencontrés.

    \item[Bartiméus :] Chat roux de la caravane, indépendant et imprévisible,
    présent durant une grande partie de l'enfance de Kael et Kamatsu.
\end{description}

\section*{Objets importants}

\begin{description}
    \item[La dent :] Canine de la première créature vaincue par Kael, portée
    directement à travers son oreille gauche depuis l'enfance.

    \item[Le pendentif :] Pierre sombre gravée d'un ancien symbole orc,
    autrefois portée par sa mère et transmise à Kael lors de sa fuite.

    \item[La dague :] Cadeau de Kamatsu au moment où leurs routes se sont
    séparées. Kael la porte comme arme autant que comme souvenir de son frère.
\end{description}